% eksperymentalne tłumaczenie na włoski

%

\section{Twierdzenie Pitagorasa} % lang-pl
\section{Teorema di Pitagora} % lang-it


Najważniejszym twierdzeniem dotyczącym trójkątów prostokątnych jest twierdzenie Pitagorasa oraz twierdzenie do niego odwrotne. % lang-pl
Il più importante teorema riguardante i triangoli rettangoli è il teorema di Pitagora e il suo inverso. % lang-it

Piszą o nim Guzicki \cite[s. 160]{guzicki_2021}. % lang-pl
Ne scrive Guzicki \cite[p. 160]{guzicki_2021}. % lang-it

% PRZECZYTANO: https://en.wikipedia.org/wiki/Pythagorean_theorem

Do sformułowania twierdzenia Pitagorasa potrzebne jest pojęcie pola, którego z pewnych powodów nigdzie formalnie nie opisaliśmy. % lang-pl

Wystarczy nam wiedzieć, że dla każdej odpowiednio regularnej figury na płaszczyźnie $F$ istnieje liczba rzeczywista $S_F \ge 0$ (zwana polem) taka, że: % lang-pl
Per formulare il teorema di Pitagora è necessario il concetto di area, che per varie ragioni non abbiamo mai definito formalmente. % lang-it

Ci basterà sapere che per ogni figura sufficientemente regolare del piano $F$ esiste un numero reale $S_F \ge 0$ (detto area) tale che: % lang-it
\begin{itemize}
    \item pole prostokąta o bokach długości $a$ oraz $b$ to $ab$; % lang-pl
    \item pole sumy (,,unii'') dwóch rozłącznych figur $F_1 \cup F_2$ jest sumą ich pól; % lang-pl
    \item pola przystających trójkątów są sobie równe. % lang-pl
    \item l'area di un rettangolo di lati $a$ e $b$ è uguale ad $ab$; % lang-it
    \item l'area dell'unione di due figure disgiunte $F_1 \cup F_2$ è la somma delle loro aree; % lang-it
    \item i triangoli congruenti hanno la stessa area. % lang-it
\end{itemize}

(Można poczytać o tym u Neugebauera \cite[s. 37]{neugebauer_2018}). % lang-pl

(Si può leggere qualcosa in proposito in Neugebauer \cite[p. 37]{neugebauer_2018}). % lang-it

\begin{theorem}[Pitagorasa] % lang-pl
\begin{theorem}[di Pitagora] % lang-it
\label{theorem_pythagorean}%
\index{twierdzenie!Pitagorasa} % lang-pl
    Dany jest trójkąt prostokątny $\triangle ABC$, w którym kąt przy wierzchołku $C$ jest prosty. % lang-pl
\index{teorema!di Pitagora}% % lang-it
    
    Na jego bokach zbudowano trzy kwadraty. % lang-pl
    Sia dato un triangolo rettangolo $\triangle ABC$, nel quale l'angolo al vertice $C$ è retto. % lang-it
    
    Sui suoi lati sono costruiti tre quadrati. % lang-it
    \begin{center}
\begin{comment}
...
\end{comment}
    \end{center}
    Suma pól jasnych kwadratów jest równa polu ciemnego kwadratu: % lang-pl
    
    La somma delle aree dei quadrati chiari è uguale all'area del quadrato scuro: % lang-it
    \begin{equation}
        |AC|^2 + |BC|^2 = |AB|^2.
    \end{equation}
    Odwrotnie, jeśli $ABC$ jest trójkątem takim, że $|AC|^2 + |BC|^2 = |AB|^2$, to trójkąt ten jest prostokątny, zaś kąt przy wierzchołku $C$ jest prosty. % lang-pl
    
    Viceversa, se $ABC$ è un triangolo tale che $|AC|^2 + |BC|^2 = |AB|^2$, allora il triangolo è rettangolo e l'angolo al vertice $C$ è retto. % lang-it
\end{theorem}

Powyższe twierdzenie przypiszemy kiedyś Pitagorasowi z Samos, choć nie wiemy dokładnie, kto i kiedy odkryje je jako pierwszy. % lang-pl

In seguito attribuiremo questo teorema a Pitagora di Samo, sebbene non sappiamo con precisione chi lo abbia scoperto per primo né quando. % lang-it
\index[persons]{Pitagoras z Samos}%
Będzie powszechnie stosowane w okresie Starego Babilonu (XX-XVI wiek p.n.e.), a więc na długo przed narodzinami Pitagorasa; pojawi się też w indyjskich i chińskich tekstach matematycznych. % lang-pl

Papirus Berlin 6619 spisany ok. 1800 roku p.n.e. na terenach państwa egipskiego zawrze zadanie, którego rozwiązaniem jest trójka $(6, 8, 10)$. % lang-pl
Esso sarà ampiamente utilizzato nel periodo dell'Antica Babilonia (XX-XVI secolo a.C.), molto prima della nascita di Pitagora; comparirà inoltre in testi matematici indiani e cinesi. % lang-it

\index{papirus Berlin 6619} % lang-pl
Jest jeszcze babilońska tabliczka Plimpton 322, także spisana ok. 1800 roku p.n.e., gdzie pojawia się trójka % lang-pl
Il papiro Berlin 6619, redatto intorno al 1800 a.C. in territorio egiziano, conterrà un problema la cui soluzione è la terna $(6, 8, 10)$. % lang-it

\index{papiro Berlin 6619}% % lang-it
Vi è inoltre la tavoletta babilonese Plimpton 322, anch'essa datata intorno al 1800 a.C., nella quale compare la terna % lang-it
\begin{equation}
    12709^2 + 13500^2 = 18541^2,
\end{equation}
co sugeruje, że jej autor znał pewną systematyczną metodę. % lang-pl

\index{tabliczka Plimpton 322} % lang-pl
il che suggerisce che il suo autore conoscesse un metodo sistematico. % lang-it

Być może twierdzenie Pitagorasa ma więcej znanych dowodów niż jakiekolwiek inne (poza prawem wzajemności reszt kwadratowych). % lang-pl
Będzie ich tak bardzo bez liku, że nie wiadomo, ile dokładnie. % lang-pl
Niektóre opierają się na rozcięciu pewnej układanki na fragmenty, przestawieniu ich i zbudowaniu innego kształtu. % lang-pl
Inne korzystają z podobieństwa trójkątów. % lang-pl
Dowód Euklidesa urzekł nas tak bardzo swoją pomysłowością, że będzie jedynym, jaki przedstawimy w całej książce! % lang-pl
\index{zasada!Cavalieriego} % lang-pl
Patrz też: Neugebauer \cite[s. 39-41]{neugebauer_2018}. % lang-pl
\index{tavoletta Plimpton 322}% % lang-it


Forse il teorema di Pitagora possiede più dimostrazioni conosciute di qualsiasi altro teorema (a eccezione della legge di reciprocità quadratica). % lang-it
Saranno così numerose da rendere impossibile stabilirne il numero esatto. % lang-it
Alcune si basano sul taglio di una figura in pezzi, sul loro riordinamento e sulla costruzione di una nuova figura. % lang-it
Altre utilizzano la similitudine dei triangoli. % lang-it
La dimostrazione di Euclide ci ha affascinato così tanto per la sua ingegnosità che sarà l'unica presentata in tutto il libro! % lang-it
\index{principio!di Cavalieri} % lang-it
Si veda anche Neugebauer \cite[p. 39--41]{neugebauer_2018}. % lang-it
\begin{proof}
    Niech $\triangle ABC$ będzie trójkątem prostokątnym, z kątem prostym przy wierzchołku $C$. % lang-pl
    
    Na bokach $BC$, $AB$, $CA$ kreślimy kolejno kwadraty $BCDE$, $ABFG$, $ACHI$ (konstrukcja kwadratu Euklidesa korzysta z postulatu równoległości). % lang-pl
    Sia $\triangle ABC$ un triangolo rettangolo con angolo retto nel vertice $C$. % lang-it
    
    Sui lati $BC$, $AB$, $CA$ costruiamo rispettivamente i quadrati $BCDE$, $ABFG$, $ACHI$ (la costruzione euclidea del quadrato utilizza il postulato delle parallele). % lang-it

    \begin{center}
\begin{comment}
...
\end{comment}
    \end{center}

    Z punktu $C$ opuszczamy wysokość na przeciwprostokątną $AB$ i przedłużamy tak, by przecięła kwadrat $ABFG$ w punktach $K$ i $L$. % lang-pl
    
    Łączymy punkty $B$ i $I$ oraz $C$ i $G$. % lang-pl
    Dal punto $C$ tracciamo l'altezza all'ipotenusa $AB$ e la prolunghiamo in modo che intersechi il quadrato $ABFG$ nei punti $K$ e $L$. % lang-it
    
    Otrzymane trójkąty $\triangle BAI$ oraz $\triangle GAC$ są przystające na mocy cechy bok-kąt-bok. % lang-pl
    Congiungiamo i punti $B$ e $I$, nonché $C$ e $G$. % lang-it
    
    Niebieski prostokąt $AGLK$ (odpowiednio: kwadrat $ACHI$) ma dwukrotnie większe pole niż trójkąt $\triangle GAC$ (trójkąt $\triangle BAI$). % lang-pl
    I triangoli ottenuti $\triangle BAI$ e $\triangle GAC$ sono congruenti per il criterio lato-angolo-lato. % lang-it
    
    (To jest zamaskowana zasada Cavalieriego!). % lang-pl
    Il rettangolo blu $AGLK$ (rispettivamente il quadrato $ACHI$) ha area doppia rispetto al triangolo $\triangle GAC$ (rispettivamente al triangolo $\triangle BAI$). % lang-it
    
    \index{zasada Cavalieriego} % lang-pl
    Zatem prostokąt $AGLK$ i kwadrat $ACHI$ mają równe pola. % lang-pl
    (Questo è un principio di Cavalieri mascherato!). % lang-it
    
    \index{principio di Cavalieri}% % lang-it
    Pertanto il rettangolo $AGLK$ e il quadrato $ACHI$ hanno la stessa area. % lang-it

    Analogicznie pokazujemy, że czerwony prostokąt $BFLK$ i kwadrat $BCDE$ mają równe pola. % lang-pl
    
    Dodajemy dwie równości stronami i otrzymujemy, że suma pól kwadratów $ACIH$ oraz $BCDE$ jest równa polu kwadratu $ABFG$. % lang-pl
    In modo analogo si dimostra che il rettangolo rosso $BFLK$ e il quadrato $BCDE$ hanno la stessa area. % lang-it
    
    Sommando membro a membro le due uguaglianze otteniamo che la somma delle aree dei quadrati $ACIH$ e $BCDE$ è uguale all'area del quadrato $ABFG$. % lang-it
\end{proof}

Według legendy Hippazos z Metapontu odkryje, że przekątna kwadratu (albo, według innej legendy, pięciokąta) nie jest współmierna z jego bokiem. % lang-pl

Secondo la leggenda, Ippaso di Metaponto scoprirà che la diagonale del quadrato (oppure, secondo un'altra versione, del pentagono) non è commensurabile con il suo lato. % lang-it
\index[persons]{Hippazos z Metapontu}%
Niewymierność liczb $\sqrt{2}$ (albo $(1 + \sqrt 5) /2$) zrujnuje pogląd szkoły pitagorejskiej, że świat opiera się na liczbach (...) i doprowadzi do utopienia Hippazosa. % lang-pl

\index{utopienie} % lang-pl
Ale jego śmierć nie cofnie rozłamu, jaki powstanie w szkole. % lang-pl
L'irrazionalità di $\sqrt{2}$ (oppure di $(1+\sqrt5)/2$) distruggerà la concezione pitagorica secondo cui il mondo si fonda sui numeri (...) e porterà all'annegamento di Ippaso. % lang-it
\index{annegamento}% % lang-it
Ma la sua morte non porrà fine alla frattura che si produrrà nella scuola. % lang-it


(Być może w tym miejscu warto dowiedzieć się o cegle Eulera.) % lang-pl
(Forse questo è il momento giusto per conoscere il mattone di Eulero.) % lang-it
\begin{corollary} % lang-it
    La lunghezza della diagonale di un rettangolo di lati $a$ e $b$ è $c = \sqrt{a^2+b^2}$. % lang-it
\end{corollary} % lang-it


\begin{corollary} % lang-pl
    Długość przekątnej prostokąta o bokach długości $a$ i $b$ wynosi $c = \sqrt{a^2 + b^2}$. % lang-pl
\end{corollary} % lang-pl
Kto umie rozwijać funkcje w szeregi potęgowe może wpaść na pomysł, by zapisać % lang-pl
Chi sa sviluppare le funzioni in serie di potenze potrebbe scrivere % lang-it
\begin{equation}
    c = a + \frac{b^2}{2a} - \frac{b^4}{8a^3} + \frac{b^6}{16a^5} - \frac{5b^8}{127a^7} + \ldots,
\end{equation}
co po obcięciu do pierwszych dwóch wyrazów daje przybliżenie znane Babilończykom (\cite[s. 6]{eves1_1972}): % lang-pl

che, troncata ai primi due termini, fornisce l'approssimazione già nota ai Babilonesi (\cite[p. 6]{eves1_1972}): % lang-it
\begin{equation}
    c \approx a + \frac{b^2}{2a}.
\end{equation}

Metoda ta daje największy błąd względny, gdy $a = b$ i wynosi on $6.06601\ldots \%$. % lang-pl
Questo metodo produce il massimo errore relativo quando $a=b$, e tale errore è pari a $6.06601\ldots \%$. % lang-it
\begin{problem} % lang-it
    Una città è circondata da una cinta muraria circolare di diametro sconosciuto, con quattro porte. % lang-it
    Un albero si trova tre nibli a nord della porta settentrionale. % lang-it
    Se ci spostiamo di nove nibli a est della porta meridionale, l'albero diventa visibile. % lang-it
    Determinare il diametro della città. % lang-it
\end{problem} % lang-it


\begin{problem} % lang-pl
    Miasto jest otoczone okrągłym murem o nieznanej średnicy, z czterema bramami. % lang-pl
    Drzewo leży trzy nible na północ od północnej bramy. % lang-pl
    Jeśli idziemy dziewięć nibli na wschód od południowej bramy, drzewo staje się widoczne. % lang-pl
    Znajdź średnicę miasta. % lang-pl
\end{problem} % lang-pl
Il problema proviene dal testo cinese \emph{Shu shu jiu zhang} (``Trattato matematico in nove sezioni'') del XIII secolo. % lang-it

Problem pochodzi z chińskiego tekstu \emph{Shu shu jiu zhang} (,,Traktat matematyczny w dziewięciu sekcjach'') z XIII wieku. % lang-pl
Jeżeli oznaczymy średnicę miasta przez $x$, odległość drzewa od miasta przez $b$, odległość do przejścia przez $a$, zaś kąt środek miasta -- punkt widoczności -- południowa brama przez $\alpha$, to dostaniemy równanie % lang-pl
Se indichiamo con $x$ il diametro della città, con $b$ la distanza dell'albero dalla città, con $a$ la distanza percorsa e con $\alpha$ l'angolo formato dal centro della città, dal punto di osservazione e dalla porta meridionale, otteniamo l'equazione % lang-it
\begin{equation}
    \tan 2 \alpha = 2 \tan \alpha + \frac ba,
\end{equation}
co po prostych przekształceniach i przyjęciu $y = \tan \alpha$ daje % lang-pl

che, dopo semplici trasformazioni e ponendo $y=\tan\alpha$, diventa % lang-it
$
    2ay^3 + by^2 - b = 0
$,
albo (pamiętając, że $y = x/2a$) % lang-pl

oppure (ricordando che $y=x/2a$) % lang-it
$
    x^3 + bx^2 - 4a^2 b = 0
$.
To równanie ma jedno rozwiązanie rzeczywiste, $x = 9$. % lang-pl
Questa equazione possiede un'unica soluzione reale, $x=9$. % lang-it


Względnie pierwsze liczby naturalne $a, b, c$ takie, że $a^2 + b^2 = c^2$ nazywamy (pierwotną) trójką pitagorejską. % lang-pl
Chiamiamo terna pitagorica (primitiva) tre numeri naturali coprimi $a,b,c$ tali che $a^2+b^2=c^2$. % lang-it

\index{trójka pitagorejska} % lang-pl
Każdą taką trójkę można otrzymać biorąc względnie pierwsze liczby $m, n$ różnej parzystości takie, że $m > n$ i kładąc $a = m^2 - n^2$, $b = 2 mn$, $c = m^2 + n^2$. % lang-pl
\index{terna pitagorica}% % lang-it

Najmniejszą taką trójką jest $(3, 4, 5)$; inna legenda (nie było w niej ani smoków, ani Hippazosa) głosi, że Egipcjanie używali tego trójkąta do wyznaczania kątów prostych u podstawy piramid. % lang-pl
Ogni tale terna può essere ottenuta scegliendo numeri coprimi $m$, $n$ di parità diversa con $m>n$ e ponendo $a=m^2-n^2$, $b=2mn$, $c=m^2+n^2$. % lang-it

La più piccola di tali terne è $(3,4,5)$; un'altra leggenda (nella quale non compaiono né draghi né Ippaso) racconta che gli Egizi usassero questo triangolo per tracciare angoli retti alla base delle piramidi. % lang-it

\begin{proposition}
    Niech $\triangle ABC$ będzie trójkątem prostokątnym, z kątem prostym przy wierzchołku $C$: % lang-pl
    
    Sia $\triangle ABC$ un triangolo rettangolo con angolo retto nel vertice $C$: % lang-it
    \begin{center}
\begin{comment}
...
\end{comment}
    \end{center}
    Mają wtedy miejsce następujące równości: % lang-pl
    
    Valgono allora le seguenti uguaglianze: % lang-it
    \begin{equation}
        h = \frac{ab}{c}, \quad
        p = \frac{b^2}{c}, \quad
        q = \frac{a^2}{c}, \quad
        h^2 = pq.
    \end{equation}
\end{proposition}

Twierdzenie Pitagorasa znajduje zastosowanie także przy wyznaczaniu niektórych miejsc geometrycznych. % lang-pl

Il teorema di Pitagora trova applicazione anche nella determinazione di alcuni luoghi geometrici. % lang-it

\begin{proposition}
    Dane są dwa różne punkty $A$ i $B$ na płaszczyźnie oraz liczba rzeczywista $c$ taka, że $2c > |AB|^2$. % lang-pl
    
    Miejscem geometrycznym punktów $P$ o własności $|AP|^2 + |BP|^2 = c$ jest okrąg o środku w środku odcinka $AB$ i średnicy $r = \sqrt{2c - |AB|^2}$. % lang-pl
    Siano dati due punti distinti $A$ e $B$ del piano e un numero reale $c$ tale che $2c>|AB|^2$. % lang-it
    
    Il luogo geometrico dei punti $P$ che soddisfano $|AP|^2+|BP|^2=c$ è una circonferenza avente centro nel punto medio di $AB$ e diametro $r=\sqrt{2c-|AB|^2}$. % lang-it
\end{proposition}

\begin{proposition}
    Dane są dwa różne punkty $A$ i $B$ na płaszczyźnie oraz liczba rzeczywista $c$. % lang-pl
    
    Wtedy miejscem geometrycznym punktów $P$ o własności $|AP|^2 - |BP|^2 = c$ jest prosta prostopadła do prostej $AB$. % lang-pl
    Siano dati due punti distinti $A$ e $B$ del piano e un numero reale $c$. % lang-it
    
    Allora il luogo geometrico dei punti $P$ tali che $|AP|^2-|BP|^2=c$ è una retta perpendicolare alla retta $AB$. % lang-it
\end{proposition}

Patrz Guzicki \cite[s. 170-173]{guzicki_2021} (Guzicki wprowadza potem osie i środki potęgowe jak w fakcie \ref{guzicki_6_11}, a następnie twierdzenie \ref{guzicki_6_13} (Carnota)). % lang-pl
Si veda Guzicki \cite[p. 170--173]{guzicki_2021} (Guzicki introduce successivamente gli assi e i centri radicali come nel fatto \ref{guzicki_6_11}, e poi il teorema \ref{guzicki_6_13} (di Carnot)). % lang-it


Spirala Teodor(os)a z Cyreny składa się z trójkątów prostokątnych stykających się ze sobą wzdłuż boków. % lang-pl
La spirale di Teodoro di Cirene è composta da triangoli rettangoli adiacenti lungo i lati. % lang-it

\index{spirala Teodor(os)a} % lang-pl
\index{spirale di Teodoro}% % lang-it
\index[persons]{Teodor(os) z Cyreny}%
Zaczynamy od trójkąta prostokątnego równoramiennego, o bokach długości $1$, $1$, $\sqrt{2}$, by następnie kreślić jednostkowy odcinek prostopadły do końca przeciwprostokątnej, łączymy go z początkiem i powtarzamy (Teodoros zatrzyma się na trójkącie, którego najdłuższy bok to $\sqrt{17}$, ponieważ następny naszedłby na ten, od którego zaczynaliśmy). % lang-pl

Si parte da un triangolo rettangolo isoscele con lati di lunghezza $1$, $1$, $\sqrt2$; quindi si costruisce un segmento unitario perpendicolare all'estremità dell'ipotenusa, lo si collega all'origine e si ripete il procedimento (Teodoro si fermò al triangolo il cui lato più lungo misura $\sqrt{17}$, poiché il successivo si sarebbe sovrapposto a quello iniziale). % lang-it
\begin{figure}[H]\centering
\begin{comment}
...
\end{comment}
\caption{Spirala Teodor(os)a, zwany czasami ślimakiem pitagorejskim} % lang-pl
\caption{Spirale di Teodoro, talvolta chiamata chiocciola pitagorica} % lang-it
\end{figure}
\index{ślimak pitagorejski} % lang-pl
\index{chiocciola pitagorica}% % lang-it

% eksperymentalne tłumaczenie na włoski

%

\index{wzór!Herona} % lang-pl

\subsection{Pole trójkąta, wzór Herona} % lang-pl
\index{formula!di Erone} % lang-it
\subsection{Area del triangolo, formula di Erone} % lang-it


W tej podsekcji znajdziemy różne wzory na pole trójkąta. % lang-pl
In questa sottosezione presenteremo diverse formule per l'area di un triangolo. % lang-it

\begin{proposition}
    Pole trójkąta o podstawie długości $a$ oraz wysokości opuszczonej na tę podstawę długości $h$ wyraża się wzorem % lang-pl
    
    L'area di un triangolo con base di lunghezza $a$ e altezza relativa a tale base di lunghezza $h$ è data dalla formula % lang-it
    \begin{equation}
        S_\triangle = \frac 1 2 a h.
    \end{equation}
    Jeśli kąt naprzeciw podstawy ma miarę $\alpha$, a ramiona długości $b, c$, to zachodzi też % lang-pl
    
    Se l'angolo opposto alla base misura $\alpha$ e gli altri due lati hanno lunghezze $b$ e $c$, allora vale anche % lang-it
    \begin{equation}
        S_\triangle = \frac 1 2 bc \sin \alpha.
    \end{equation}
\end{proposition}

Pisze o tym Guzicki \cite[s. 255]{guzicki_2021}. % lang-pl

Se ne occupa Guzicki \cite[p.~255]{guzicki_2021}. % lang-it

\begin{proposition}[wzór Herona, 60 r.n.e.?] % lang-pl
\begin{proposition}[formula di Erone, 60 a.C.?] % lang-it
\label{prp_heron}%
    Pole trójkąta o bokach długości $a, b, c$ wyraża się wzorem % lang-pl
    
    L'area di un triangolo con lati di lunghezza $a$, $b$, $c$ è data dalla formula % lang-it
    \begin{equation}
        S = \sqrt{p(p-a)(p-b)(p-c)},
    \end{equation}
    gdzie $p = \frac 1 2 (a + b + c)$ oznacza połowę obwodu. % lang-pl
    
    dove $p = \frac 1 2 (a + b + c)$ denota il semiperimetro. % lang-it
\end{proposition}

Wynik przypiszemy Heronowi z Aleksandrii, poda dość skomplikowany dowód w~dziele \emph{,,Metrica''}. % lang-pl
Attribuiremo questo risultato a Erone di Alessandria, che ne fornirà una dimostrazione piuttosto complessa nell'opera \emph{Metrica}. % lang-it
\index[persons]{Erone di Alessandria}% % lang-it


\index[persons]{Heron z Aleksandrii}% % lang-pl
Coxeter \cite[s. 12]{coxeter_1991} (\cite[s. 28]{coxeter_1967}) przeczyta van der Waerdena \cite[s. 228, 277]{waerden_1961} i przez to stwierdzi, że wzór był już znany Archimedesowi dwieście lat wcześniej. % lang-pl
Coxeter \cite[p.~12]{coxeter_1991} (\cite[p.~28]{coxeter_1967}) leggerà van der Waerden \cite[pp.~228, 277]{waerden_1961} e concluderà che la formula era già nota ad Archimede duecento anni prima. % lang-it

% czemu jest cytowane angielskie i polskie wydanie?

\index[persons]{Archimedes}% % lang-pl
\index[persons]{Archimede}% % lang-it

Pompe w $\Delta_{93}^9$ pokaże, jak dojść do dowodu Sidneya Kunga ,,bez rachowania''. % lang-pl

Pompe mostrerà in $\Delta_{93}^9$ come ottenere la dimostrazione di Sidney Kung ``senza fare calcoli''. % lang-it
% https://www.deltami.edu.pl/1993/09/jak-udowodnic-wzor-herona/

\index[persons]{Kung, Sidney}%

Guzicki \cite[s. 165-169]{guzicki_2021} wyprowadzi go z twierdzenia Pitagorasa i~wspomni, że prosty geometryczny dowód poda(ł) też Euler. % lang-pl
Guzicki \cite[pp.~165--169]{guzicki_2021} lo dedurrà dal teorema di Pitagora e ricorderà che una semplice dimostrazione geometrica fu data anche da Eulero. % lang-it
\index[persons]{Euler, Leonhard}% % lang-it


\index[persons]{Euler, Leonhard}% % lang-pl
\index{twierdzenie!Pitagorasa} % lang-pl
\index{teorema!di Pitagora} % lang-it


Wreszcie Bogdańska, Neugebauer \cite[s. 92]{neugebauer_2018} skorzystają z twierdzenia kosinusów. % lang-pl
Infine Bogdańska e Neugebauer \cite[p.~92]{neugebauer_2018} utilizzeranno il teorema del coseno. % lang-it


Wzór Herona pojawi się wielokrotnie w innych numerach Delty: $\Delta_{14}^3$, $\Delta_{25}^{12}$ (jako pierwsz wpis na liście Marcusa Bakera stu dziesięciu różnych przepisów na pole trójkąta). % lang-pl
La formula di Erone comparirà più volte in altri numeri di Delta: $\Delta_{14}^3$, $\Delta_{25}^{12}$ (come prima voce nell'elenco di Marcus Baker di centodieci diverse formule per l'area del triangolo). % lang-it

% https://www.deltami.edu.pl/2014/03/heron-uogolniony/
% https://www.deltami.edu.pl/2025/12/heron-heron-po-trzykroc-heron/
% https://archive.org/details/jstor-1967384/page/n5/mode/2up
% https://archive.org/details/jstor-1967393

Istnieje częściowe uogólnienie wzoru Herona dla czworokątów: % lang-pl

Esiste una generalizzazione parziale della formula di Erone ai quadrilateri: % lang-it

\begin{proposition}[wzór Brahmagupty] % lang-pl
\begin{proposition}[formula di Brahmagupta] % lang-it
\label{brahmagupta_formula}%
\index{czworokąt!cykliczny}% % lang-pl
\index{wzór!Brahmagupty}% % lang-pl
    Pole czworokąta wypukłego, na którym można opisać okrąg, o bokach długości $a, b, c, d$ wyraża się wzorem % lang-pl
\index{quadrilatero!ciclico}% % lang-it
\index{formula!di Brahmagupta}% % lang-it
    
    L'area di un quadrilatero convesso ciclico con lati di lunghezza $a$, $b$, $c$, $d$ è data dalla formula % lang-it
    \begin{equation}
        S = \sqrt{(p-a)(p-b)(p-c)(p-d)},
    \end{equation}
    gdzie $p$ oznacza połowę obwodu. % lang-pl
    
    dove $p$ denota il semiperimetro. % lang-it
\end{proposition}

Dowód polega na zastosowaniu wzoru Herona dwa razy do trójkątów podobnych lub ponownym skorzystaniu z dobrodziejstw trygonometrii. % lang-pl
La dimostrazione consiste nell'applicare due volte la formula di Erone a opportuni triangoli oppure nel ricorrere nuovamente alla trigonometria. % lang-it


Założenie o wypukłości i cykliczności można pominąć, jak odkryją równocześnie Carl Bretschneider i Karl von Staudt: % lang-pl
Le ipotesi di convessità e ciclicità possono essere eliminate, come scopriranno indipendentemente Carl Bretschneider e Karl von Staudt: % lang-it

\index[persons]{Staudt@von Staudt, Karl}%
\index[persons]{Bretschneider, Carl}%

\begin{proposition}[wzór Bretschneidera, 1842] % lang-pl
\index{wzór!Bretschneidera}% % lang-pl
    Pole czworokąta o bokach długości $a, b, c, d$ i sumie miar dwóch przeciwległych kątów wewnętrznych $2 \varphi$ wyraża się wzorem % lang-pl
\begin{proposition}[formula di Bretschneider, 1842] % lang-it
\index{formula!di Bretschneider}% % lang-it
    
    L'area di un quadrilatero con lati di lunghezza $a$, $b$, $c$, $d$ e somma di due angoli opposti pari a $2\varphi$ è data dalla formula % lang-it
    \begin{equation}
        S = \sqrt{(p-a)(p-b)(p-c)(p-d) - abcd \cos^2  \varphi},
    \end{equation}
    gdzie $p$ oznacza połowę obwodu. % lang-pl
    
    dove $p$ denota il semiperimetro. % lang-it
\end{proposition}

Inny, równie przydatny wzór znajdzie Julian Coolidge \cite{coolidge_1939}: % lang-pl

Julian Coolidge \cite{coolidge_1939} troverà un'altra formula altrettanto utile: % lang-it

\index[persons]{Coolidge, Julian}%

pole czworokąta o bokach długości $a, b, c, d$ i przekątnych długości $e, f$ wyraża się wzorem % lang-pl

l'area di un quadrilatero con lati di lunghezza $a$, $b$, $c$, $d$ e diagonali di lunghezza $e$, $f$ è data dalla formula % lang-it

\begin{equation}
    S = \sqrt{(p-a)(p-b)(p-c)(p-d) - \frac 1 4 (ac + bd + pq)(ac + bd - pq)},
\end{equation}

gdzie $p$ oznacza raz jeszcze (i nie ostatni) połowę obwodu. % lang-pl

dove $p$ denota ancora una volta (e non per l'ultima) il semiperimetro. % lang-it

% % 2.4.x, nie 2.5

%